\section{Conclusion}\label{sec:conclusion}
Our paper makes three main contributions.
First,
we analytically show that distance-3 (-5) T-state cultivation
has fault distance 2 (3) before its escape stage;
this is the main reason for the T--S discrepancy
in logical error rate.
Second,
we derive a general formula for,
and an algorithm for calculating,
the acceptance probability of
an arbitrary encoded mixed state suffering a physical Clifford error.
Third,
we design 2D-local flags that restore
the degraded fault distance of T-state cultivation,
and which are practical for distance 5.
These flags demonstrate the utility of considering
coherent non-Pauli errors arising from non-Clifford circuits
for improving their fault tolerance.
There are various avenues for future work.

Surti et al.~\cite{Surti2026} show that
in many logical magic state preparation circuits
(even those with multiple layers of non-Cliffords),
each Pauli error propagates to a Clifford error.
\Cref{prop:resulting_fidelity} could be generalised to
the same level of generality as \cref{theorem:probability_trivial_syndrome}
using the same proof technique.
These two facts could then be combined with
\cref{alg:probability_trivial_syndrome}
to quickly and analytically verify the fault distance
(by enumerating low-fault-count fault configurations)
of other circuits, such as
the double checks~\cite{Chen2026,Vaknin2025,Claes2025}
or the constant-depth logical measurement by gauging~\cite{Hetenyi2026}
from more recent cultivation proposals.
The transversal-T code-switching protocol
from Daguerre et al.~\cite{Daguerre2025,Daguerre2025a},
which is an alternative to cultivation,
could also be analysed in this way.

One motivation for designing the Z-flagged circuits
was to retain the use of the double check
for its low spacetime cost under 2D connectivity constraints.
As \cref{sec:z_flagged_double_checks} mentions however,
the Z flags themselves complicate the connectivity
to next-nearest-neighbour;
it would be interesting to devise a qubit layout
that maintains nearest-neighbour connectivity.
This should be combined with
a more rigorous and systematic flag design process
than the hand-optimised one described in \cref{sec:flag_design_process}.
Additionally,
it is unknown how degraded the fault distance is for higher-distance MSC;
if the trend continues from MSC-3 and -5,
our Z flags could become increasingly effective at improving
the logical error rate.